\documentclass[12pt]{article}
\usepackage[utf8]{inputenc}
\usepackage[T1]{fontenc}
\usepackage[margin=1in]{geometry}
\usepackage{setspace}
\usepackage{xcolor}
\usepackage[colorlinks=true, urlcolor=blue!60!black, linkcolor=blue!60!black]{hyperref}
\setlength{\parindent}{2em}
\setlength{\parskip}{1.25em}
\renewcommand{\baselinestretch}{1.0}
\begin{document}
\begin{center}
{\LARGE \textit{Clinton v. City of New York}, 524 U.S. 417 (1998)\par}
\vspace{0.8em}
{\large Case Brief --- Four-Page Edition\par}
\end{center}
\vspace{1em}

\section{Memory Jogger}

The Court held 6--3 that the Line Item Veto Act violates the Presentment Clause: when the President canceled a Medicaid spending provision and a targeted tax benefit after signing them into law, he in effect repealed parts of duly enacted statutes without bicameral passage and presentment (\href{https://supreme.justia.com/cases/federal/us/524/417/}{524 U.S. 417}).

\section{Detailed Case Facts}

The Line Item Veto Act of 1996 let the President, within five days after signing a bill, ``cancel in whole'' any dollar amount of discretionary budget authority, any item of new direct spending, or any limited tax benefit, if he determined that cancellation would ``(i) reduce the Federal budget deficit; (ii) not impair any essential Government functions; and (iii) not harm the national interest.'' A canceled provision lost ``legal force or effect'' unless Congress passed a disapproval bill, itself subject to veto, and a ``lockbox'' devoted the savings to deficit reduction.

After a suit by Members of Congress was dismissed for lack of standing in \textit{Raines v. Byrd} (1997), President Clinton on August 11, 1997, canceled \S~4722(c) of the Balanced Budget Act, which had waived the federal right to recoup up to \$2.6 billion in Medicaid provider taxes from New York, and \S~968 of the Taxpayer Relief Act, which let owners of food processing facilities defer capital gains when selling to farmers' cooperatives.

\section{Procedural History}

The City of New York, its hospital corporation, hospital associations, and unions sued in the District of Columbia, as did Snake River Potato Growers, an Idaho cooperative whose purchase negotiations ended after the cancellation, and its vice-chairman, Mike Cranney. The district court consolidated the suits, found standing, and held the Act unconstitutional. The Supreme Court heard expedited direct appeals on April 27 and affirmed on June 25, 1998.

\section{Judicial Votes}

Affirmed, \textbf{6--3}.
\begin{itemize}
\item \textbf{Majority:} Justice Stevens, joined by Chief Justice Rehnquist and Justices Kennedy, Souter, Thomas, and Ginsburg.
\item \textbf{Concurrence:} Justice Kennedy.
\item \textbf{Concurring in part and dissenting in part:} Justice Scalia, joined by Justice O'Connor, and by Justice Breyer as to Part III.
\item \textbf{Dissent:} Justice Breyer, joined by Justices O'Connor and Scalia as to Part III.
\end{itemize}

\section{Holding}

Both sets of appellees had statutory and Article III standing. The Act's cancellation procedures violate the Presentment Clause, Article I, \S~7, cl.~2, because they authorize the President to amend or repeal portions of duly enacted statutes without following the exclusive lawmaking procedures of Article I, \S~7. Any such new procedure must come through an Article V amendment.

\section{Analysis of Opinions}

\textbf{Majority (Stevens).} Reading ``individual'' to mean ``person,'' Stevens distinguished \textit{Raines}: canceling \S~4722(c) revived a multibillion-dollar contingent liability for New York, and canceling \S~968 stripped Snake River of a ``bargaining chip.'' On the merits, following \textit{INS v. Chadha}, cancellation was in legal and practical effect a partial repeal, and repeal must satisfy Article I, \S~7. A constitutional return comes before a bill becomes law and covers the whole bill; cancellation comes after and strikes only a part. Constitutional silence is ``equivalent to an express prohibition,'' as George Washington understood that a President must ``approve all the parts of a Bill, or reject it in toto.'' \textit{Field v. Clark} differed because its tariff suspension depended on later foreign conditions, was mandatory, and executed Congress's policy; traditional ``sums not exceeding'' appropriations never let the President ``change the text of duly enacted statutes.''

\textbf{Kennedy (concurring).} Answering Breyer, Kennedy explained that separation of powers protects liberty both between branches and between government and citizens. Letting the President target states or taxpayers undermines political liberty, and ``one Congress cannot yield up its own powers, much less those of other Congresses to follow.''

\textbf{Scalia (concurring in part, dissenting in part).} Snake River lacked standing. As to New York, the Balanced Budget Act was duly enacted before any cancellation, so the real issue was delegation; from the First Congress to \textit{Train v. City of New York}, Congress has let Presidents ``decline to spend.'' The Act's title ``succeeded in faking out the Supreme Court.''

\textbf{Breyer (dissenting).} Canceling under a preexisting statutory proviso executes the statute rather than repealing it. Given a \$1.5 trillion budget, the Act did not encroach on Congress and satisfied the ``intelligible principle'' test.

\section{Examples: Future Applications}

\begin{enumerate}
\item \textbf{Same-Side (Invalid): Earmark Cancellation.} A statute lets the President cancel earmarks within ten days of signing. It repeats the defect: action after enactment, on part of a law, on the same facts Congress weighed.
\item \textbf{Same-Side (Invalid): Treasury Tax-Credit Veto.} The Secretary may void any newly enacted tax credit benefiting fewer than one hundred taxpayers as poor policy, failing all three \textit{Field v. Clark} distinctions.
\item \textbf{Opposite-Side (Valid): Expedited Rescission.} Presidential rescission packages get fast-track up-or-down votes but take effect only if both Houses pass them and the President signs.
\item \textbf{Opposite-Side (Valid): Contingent Tariffs.} The President must raise duties on any country later found to discriminate against American goods, as in \textit{Field v. Clark}.
\item \textbf{Fence-Sitter: Non-Expenditure Directives.} A statute deems all appropriations ``sums not exceeding'' and lets the President bar spending on line items to ``promote fiscal discipline.'' The majority conceded such appropriations are valid, yet permanently blocking an earmark resembles cancellation.
\end{enumerate}

\section{Critique}

\textbf{Formalism versus substance.} Following Scalia, critics ask whether \textit{Clinton} elevates drafting magic words over substance: if Congress may appropriate ``sum[s] not exceeding'' an amount and leave allocation to the President, as Stevens acknowledged, the formal line between ``declining to spend'' and ``erasing statutory text'' may merely force Congress to draft permissive appropriations rather than constrain executive power.

\textbf{Severability and taxes.} Stevens declined to reach severability, and Scalia's defense of spending discretion did not clearly extend to canceling \textit{tax} provisions, which he avoided through standing, leaving open whether revoking a tax benefit is more clearly legislative than declining to spend.

\section{Key Quotations}

\begin{quote}
``Failure of political will does not justify unconstitutional remedies\ldots\ Liberty is always at stake when one or more of the branches seek to transgress the separation of powers\ldots\ Abdication of responsibility is not part of the constitutional design.''
\end{quote}
\noindent --- Justice Kennedy, concurring, 524 U.S. at 449, 450, 452.

\begin{quote}
``If the Line Item Veto Act were valid, it would authorize the President to create a different law---one whose text was not voted on by either House of Congress or presented to the President for signature\ldots\ If there is to be a new procedure in which the President will play a different role in determining the final text of what may `become a law,' such change must come not by legislation but through the amendment procedures set forth in Article V of the Constitution.''
\end{quote}
\noindent --- Justice Stevens, for the Court, 524 U.S. at 448--449.

\begin{quote}
``The constitutional return takes place before the bill becomes law; the statutory cancellation occurs after the bill becomes law. The constitutional return is of the entire bill; the statutory cancellation is of only a part.''
\end{quote}
\noindent --- Justice Stevens, for the Court, 524 U.S. at 439.

\end{document}
